\documentclass[10pt]{beamer}
\usetheme{Madrid}
\usepackage{booktabs}
\usepackage{graphicx}
\usepackage[T1]{fontenc}
\setbeamertemplate{navigation symbols}{}

% ======================= EDIT YOUR DETAILS =======================
\newcommand{\studentname}{Aflah Muhammed P}
% Change the university roll/registration number:
\newcommand{\rollno}{TVE23CSXXX}
% Change your class roll number:
\newcommand{\rollclass}{Roll No: XX}
% Change your guide name:
\newcommand{\guidename}{Prof. Guide Name}
% Change department / college / short-institute if needed:
\newcommand{\deptname}{Department of Computer Science and Engineering}
\newcommand{\collegename}{College of Engineering Trivandrum}
\newcommand{\shortinst}{CET}
% Change the seminar date:
\newcommand{\seminardate}{\today}
% =================================================================

\title[B.Tech Seminar]{When Helpful Agents Overshare: How Simple Hidden Prompts Leak Personal Data}
\subtitle{Based on: Simple Prompt Injection Attacks Can Leak Personal Data Observed by LLM Agents During Task Execution (arXiv:2506.01055, 2025)}
\author[\studentname\ (\shortinst)]{\studentname \\ \small \rollno \\ \small \rollclass \\ \small Guide: \guidename}
\institute[]{\deptname \\ \collegename}
\date{\seminardate}

\begin{document}

\begin{frame}[plain]
  \titlepage
\end{frame}

\begin{frame}{Table of Contents}
  \tableofcontents
\end{frame}

\section{Introduction}
\begin{frame}{Introduction}
\begin{itemize}
  \item AI agents now call real tools, not just chat
  \item Agents must read outside, possibly untrusted content
  \item Hidden instructions can hijack what agents do
  \item Agents see personal data while doing normal tasks
  \item This work measures how easily that data leaks
\end{itemize}
\end{frame}

\section{Literature Survey}
\begin{frame}{Literature Survey / Background}
  \small
  \begin{tabular}{@{}p{2.7cm}p{2.3cm}p{5.6cm}@{}}
    \toprule
    \textbf{Title} & \textbf{Author} & \textbf{Summary} \\
    \midrule
    AgentDojo: A Dynamic Environment to Evaluate Attacks and Defenses for LLM Agents & Debenedetti et al. (2024) & Standard benchmark with realistic tool-using tasks and injection tests across banking, email, Slack, and travel; this paper extends its banking suite. \\
    \addlinespace
    Not What You've Signed Up For: Compromising LLM-Integrated Applications with Indirect Prompt Injection & Greshake et al. (2023) & Showed attackers can hide instructions in content apps retrieve, hijacking real-world LLM applications without any direct access to the user. \\
    \addlinespace
    Defeating Prompt Injections by Design (CaMeL) & Debenedetti et al. (2025) & Proposes a design-based defense separating trusted control from untrusted data flow; the paper notes it is not yet tested against these attacks. \\
    \addlinespace
    \bottomrule
  \end{tabular}
\end{frame}

\section{Motivation}
\begin{frame}{Motivation}
\begin{itemize}
  \item Agents are entering sensitive domains like banking
  \item Prior benchmarks tested generic tasks, not data leakage
  \item Personal data is observed during ordinary task execution
  \item Even simple, non-technical attacks may succeed
\end{itemize}
\end{frame}

\section{Problem Statement}
\begin{frame}{Problem Statement}
  \begin{block}{}
  Can short, plain-language instructions hidden in content that a tool-calling agent reads trick it into leaking personal data it observed while doing a legitimate task, and how well do existing defenses stop this?
  \end{block}
\end{frame}

\section{Objectives}
\begin{frame}{Objectives}
\begin{itemize}
  \item Measure data-leakage risk in a realistic banking agent
  \item Design data flow based prompt injection attacks
  \item Integrate these attacks into the AgentDojo benchmark
  \item Compare vulnerability across six different language models
  \item Evaluate whether built-in defenses prevent leakage
\end{itemize}
\end{frame}

\section{Methodology}
\begin{frame}[allowframebreaks]{Methodology / Approach}
\begin{itemize}
  \item Build a fictitious tool-calling banking agent
  \item Create a synthetic human-AI banking conversation dataset
  \item Extend AgentDojo banking suite with leakage attacks
  \item Inject an Important message style malicious instruction
  \item Vary requested data, wording, and attacker knowledge
  \item Test four AgentDojo defenses and measure trade-offs
\end{itemize}
\end{frame}

\section{Key Concepts}
\begin{frame}[allowframebreaks]{Key Concepts}
  \begin{description}
    \item[Tool-calling agent] LLM that invokes external tools to read data and take actions on a user's behalf.
    \item[Indirect prompt injection] Hidden malicious instructions placed in content the agent reads, not typed by the user.
    \item[Data flow based attack] Injection that redirects personal data the agent observed toward an attacker-controlled destination.
    \item[Attack success rate (ASR)] Share of cases where the agent actually performs the attacker's intended data leak.
    \item[Utility trade-off] Defenses lower attack success but often reduce how many real tasks succeed.
  \end{description}
\end{frame}

\section{System Architecture}
\begin{frame}{System Architecture / How It Works}
  \begin{block}{Overview}
  The pipeline starts with a user giving the banking agent a normal instruction, such as summarizing recent transactions. The agent, powered by a language model, calls banking tools that return data like statements or messages, and one of those tool outputs secretly contains a hidden injected instruction. The agent reads the combined content and decides its next action; if it obeys the injection, it calls a tool such as send email to forward the user's personal data to an attacker address. A defense layer can sit between the tool outputs and the model, for example wrapping data in delimiters, scanning it with a detector, repeating the user's request, or limiting available tools, so only clean or allowed content reaches the model.
  \end{block}
  % TIP: insert a diagram here, e.g.:
  % \begin{center}\includegraphics[width=0.7\textwidth]{your-figure.png}\end{center}
\end{frame}

\section{Results}
\begin{frame}[allowframebreaks]{Results / Findings}
\begin{itemize}
  \item About 20 percent average attack success on 16 tasks
  \item Task success drops 15 to 50 percentage points under attack
  \item Extended 48-task set shows about 15 percent average ASR
  \item Llama-4 (17B) worst near 40 percent; GPT-4o most resilient
  \item Data-extraction and authorization tasks are most vulnerable
\end{itemize}
\end{frame}

\section{Discussion}
\begin{frame}{Discussion}
\begin{itemize}
  \item Safety training blocks passwords but not all personal data
  \item Vulnerability depends on task type, not only the model
  \item Defenses trade security against completing legitimate tasks
  \item No built-in defense fully stopped leakage on extended tasks
\end{itemize}
\end{frame}

\section{Challenges}
\begin{frame}{Limitations \& Challenges}
\begin{itemize}
  \item Fictitious banking setup, not the full adversarial range
  \item Results depend on chosen benchmark, dataset, and models
  \item Synthetic conversations may miss real user behavior
  \item Design-based defenses like CaMeL not yet evaluated
\end{itemize}
\end{frame}

\section{Future Work}
\begin{frame}{Future Work}
\begin{itemize}
  \item Map which sensitive data types models inherently resist leaking
  \item Evaluate design-based defenses such as CaMeL
  \item Extend testing to insurance, stock, and cryptocurrency domains
  \item Develop and test more sophisticated injection attacks
\end{itemize}
\end{frame}

\section{Conclusion}
\begin{frame}{Conclusion}
\begin{itemize}
  \item Simple hidden prompts can make agents leak observed personal data
  \item Average attack success roughly 15 to 20 percent across banking tasks
  \item Task structure and data type strongly shape vulnerability
  \item Current defenses help but sacrifice task performance
\end{itemize}
\end{frame}

\section{References}
\begin{frame}[allowframebreaks]{References}
  \footnotesize
  \begin{enumerate}
    \item Simple Prompt Injection Attacks Can Leak Personal Data Observed by LLM Agents During Task Execution, Meysam Alizadeh, Zeynab Samei, Daria Stetsenko, Fabrizio Gilardi, arXiv:2506.01055, 2025
    \item AgentDojo: A Dynamic Environment to Evaluate Attacks and Defenses for LLM Agents, Edoardo Debenedetti, Jie Zhang, Mislav Balunovic, Luca Beurer-Kellner, Marc Fischer, Florian Tramer, arXiv:2406.13352, 2024
    \item Not What You've Signed Up For: Compromising Real-World LLM-Integrated Applications with Indirect Prompt Injection, Kai Greshake et al., ACM Workshop on AI and Security, 2023
    \item Defeating Prompt Injections by Design, Edoardo Debenedetti et al., arXiv:2503.18813, 2025
    \item Defending Against Indirect Prompt Injection Attacks With Spotlighting, Keegan Hines et al., arXiv:2403.14720, 2024
  \end{enumerate}
\end{frame}

\begin{frame}[plain]
  \begin{center}
    {\Huge Thank You}\\[1em]
    {\large Questions?}
  \end{center}
\end{frame}

\end{document}